% GENERATED FILE. Edit canonical lessons, then run npm run generate:books.
% canonical-source-hash: fnv1a64:1e11c281ae421749
% canonical-lessons: GU-C453-bivu, GU-C453-dhak, GU-C453-thatharvu, GU-C453-tras, GU-C453-bhiti

\chapter{A Sudden Fright}
\label{ch:gu-a-sudden-fright}
\hlchaptermodality{453}

\begin{chapteropening}
\textbf{By the end of this chapter:} \emph{I can say what frightened me and how.}

Complete the last lesson of chapter 453: I can say what frightened me and how.
\end{chapteropening}

\section[\texorpdfstring{\gu{બીવું} (bīvũ)}{બીવું (bīvũ)}]{\gu{બીવું} (bīvũ) — to be scared}

\label{lesson:GU-C453-bivu}



\begin{quote}
Before the new one: say the Gujarati for to seethe, then the Gujarati for strife.
\end{quote}

\subsection*{You'll want to know: \gu{બીવું}}
\textbf{\gu{બીવું}} — \emph{bīvũ} — \textquotedblleft{}to be scared\textquotedblright{}.

\begin{quote}
\gu{નાની} \gu{બહેન} \gu{અંધારામાં} \gu{બીએ} \gu{છે}.
\end{quote}

\begin{quote}
\emph{nānī bahen andhārāmā̃ bīe che.}
\end{quote}

\begin{quote}
My little sister is scared of the dark.
\end{quote}

Past (he …): \gu{તે} \gu{બીધો}, \emph{te bīdho}.

\begin{cousinweb}
\textbf{\gu{બીવું}} (\emph{bīvũ}) is inherited from Sanskrit \emph{bibhēti}, \textquotedblleft{}fears\textquotedblright{}; CDIAL lists the Gujarati verb (CDIAL 9241). The noun \textbf{\gu{બીક}} (\emph{bīk}), \textquotedblleft{}fear\textquotedblright{}, belongs to the same family.
\end{cousinweb}

\begin{grammarlens}[title={What you've built}]
One more word for describing what you lived and felt.
\end{grammarlens}

\subsection*{Your turn}
\begin{itemize}
\raggedright
  \item \emph{Say it:} \emph{bīvũ}
  \item \emph{Say it:} \emph{bīvũ}, once more
  \item \emph{Say it:} \emph{nānī bahen andhārāmā̃ bīe che.} — My little sister is scared of the dark.
  \item \emph{Recall:} say the Gujarati for to seethe, then the Gujarati for strife, then say \emph{bīvũ} again
\end{itemize}

\begin{tcolorbox}[breakable,colback=teal!4,colframe=teal!35!black,title={Before you move on}]
What does \gu{બીવું} mean? (\textquotedblleft{}To be scared\textquotedblright{}.) Say it once more.
\end{tcolorbox}
\section[\texorpdfstring{\gu{ધાક} (dhāk)}{ધાક (dhāk)}]{\gu{ધાક} (dhāk) — awe; dread of someone}

\label{lesson:GU-C453-dhak}



\begin{quote}
Before the new one: say the Gujarati for strife, then the Gujarati for to be scared.
\end{quote}

\subsection*{You'll want to know: \gu{ધાક}}
\textbf{\gu{ધાક}} — \emph{dhāk} — \textquotedblleft{}awe; dread of someone\textquotedblright{}.

\begin{quote}
\gu{ગામમાં} \gu{તે} \gu{સાહેબની} \gu{ઘણી} \gu{ધાક} \gu{હતી}.
\end{quote}

\begin{quote}
\emph{gāmmā̃ te sāhebnī ghaṇī dhāk hatī.}
\end{quote}

\begin{quote}
The whole village stood in awe of that official.
\end{quote}

A feminine noun.

\begin{cousinweb}
\textbf{\gu{ધાક}} (\emph{dhāk}) goes back to a reconstructed \emph{*dhākka-}, \textquotedblleft{}pomp, boast\textquotedblright{}; CDIAL lists the Gujarati word with the sense \textquotedblleft{}awe\textquotedblright{} (CDIAL 6769).
\end{cousinweb}

\begin{grammarlens}[title={What you've built}]
One more word for describing what you lived and felt.
\end{grammarlens}

\subsection*{Your turn}
\begin{itemize}
\raggedright
  \item \emph{Say it:} \emph{dhāk}
  \item \emph{Say it:} \emph{dhāk}, once more
  \item \emph{Say it:} \emph{gāmmā̃ te sāhebnī ghaṇī dhāk hatī.} — The whole village stood in awe of that official.
  \item \emph{Recall:} say the Gujarati for strife, then the Gujarati for to be scared, then say \emph{dhāk} again
\end{itemize}

\begin{tcolorbox}[breakable,colback=teal!4,colframe=teal!35!black,title={Before you move on}]
What does \gu{ધાક} mean? (\textquotedblleft{}Awe; dread of someone\textquotedblright{}.) Say it once more.
\end{tcolorbox}
\section[\texorpdfstring{\gu{થથરવું} (thatharvũ)}{થથરવું (thatharvũ)}]{\gu{થથરવું} (thatharvũ) — to shudder; to tremble with fear}

\label{lesson:GU-C453-thatharvu}



\begin{quote}
Before the new one: say the Gujarati for to be scared, then the Gujarati for awe.
\end{quote}

\subsection*{You'll want to know: \gu{થથરવું}}
\textbf{\gu{થથરવું}} — \emph{thatharvũ} — \textquotedblleft{}to shudder; to tremble with fear\textquotedblright{}.

\begin{quote}
\gu{સાપ} \gu{જોઈને} \gu{છોકરી} \gu{થથરી} \gu{ગઈ}.
\end{quote}

\begin{quote}
\emph{sāp joīne chokrī thathrī gaī.}
\end{quote}

\begin{quote}
The girl shuddered when she saw the snake.
\end{quote}

Past (she …): \gu{તે} \gu{થથરી} \gu{ગઈ}, \emph{te thathrī gaī}.

\begin{cousinweb}
\textbf{\gu{થથરવું}} (\emph{thatharvũ}) goes back to a reconstructed \emph{*thar-}, \textquotedblleft{}tremble\textquotedblright{}; CDIAL lists the Gujarati verb (CDIAL 6092).
\end{cousinweb}

\begin{grammarlens}[title={What you've built}]
One more word for describing what you lived and felt.
\end{grammarlens}

\subsection*{Your turn}
\begin{itemize}
\raggedright
  \item \emph{Say it:} \emph{thatharvũ}
  \item \emph{Say it:} \emph{thatharvũ}, once more
  \item \emph{Say it:} \emph{sāp joīne chokrī thathrī gaī.} — The girl shuddered when she saw the snake.
  \item \emph{Recall:} say the Gujarati for to be scared, then the Gujarati for awe, then say \emph{thatharvũ} again
\end{itemize}

\begin{tcolorbox}[breakable,colback=teal!4,colframe=teal!35!black,title={Before you move on}]
What does \gu{થથરવું} mean? (\textquotedblleft{}To shudder; to tremble with fear\textquotedblright{}.) Say it once more.
\end{tcolorbox}
\section[\texorpdfstring{\gu{ત્રાસ} (trās)}{ત્રાસ (trās)}]{\gu{ત્રાસ} (trās) — torment; harassment}

\label{lesson:GU-C453-tras}



\begin{quote}
Before the new one: say the Gujarati for awe, then the Gujarati for to shudder.
\end{quote}

\subsection*{You'll want to know: \gu{ત્રાસ}}
\textbf{\gu{ત્રાસ}} — \emph{trās} — \textquotedblleft{}torment; harassment\textquotedblright{}.

\begin{quote}
\gu{મચ્છરનો} \gu{ત્રાસ} \gu{આખી} \gu{રાત} \gu{ચાલ્યો}.
\end{quote}

\begin{quote}
\emph{maccharno trās ākhī rāt cālyo.}
\end{quote}

\begin{quote}
The mosquitoes tormented us all night.
\end{quote}

A masculine noun.

\begin{cousinweb}
\textbf{\gu{ત્રાસ}} (\emph{trās}) is inherited from Sanskrit \emph{trāsa}, \textquotedblleft{}fear\textquotedblright{}; CDIAL lists the Gujarati word, \textquotedblleft{}fear, vexation\textquotedblright{} (CDIAL 6013), and today it is mostly the vexation.
\end{cousinweb}

\begin{grammarlens}[title={What you've built}]
One more word for describing what you lived and felt.
\end{grammarlens}

\subsection*{Your turn}
\begin{itemize}
\raggedright
  \item \emph{Say it:} \emph{trās}
  \item \emph{Say it:} \emph{trās}, once more
  \item \emph{Say it:} \emph{maccharno trās ākhī rāt cālyo.} — The mosquitoes tormented us all night.
  \item \emph{Recall:} say the Gujarati for awe, then the Gujarati for to shudder, then say \emph{trās} again
\end{itemize}

\begin{tcolorbox}[breakable,colback=teal!4,colframe=teal!35!black,title={Before you move on}]
What does \gu{ત્રાસ} mean? (\textquotedblleft{}Torment; harassment\textquotedblright{}.) Say it once more.
\end{tcolorbox}
\section[\texorpdfstring{\gu{ભીતિ} (bhīti)}{ભીતિ (bhīti)}]{\gu{ભીતિ} (bhīti) — apprehension; fear of what may come}

\label{lesson:GU-C453-bhiti}



\begin{quote}
Before the new one: say the Gujarati for to shudder, then the Gujarati for torment.
\end{quote}

\subsection*{You'll want to know: \gu{ભીતિ}}
\textbf{\gu{ભીતિ}} — \emph{bhīti} — \textquotedblleft{}apprehension; fear of what may come\textquotedblright{}.

\begin{quote}
\gu{ભૂકંપ} \gu{પછી} \gu{ગામમાં} \gu{ઘણી} \gu{ભીતિ} \gu{હતી}.
\end{quote}

\begin{quote}
\emph{bhūkamp pachī gāmmā̃ ghaṇī bhīti hatī.}
\end{quote}

\begin{quote}
After the earthquake there was great apprehension in the village.
\end{quote}

A feminine noun.

\begin{cousinweb}
\textbf{\gu{ભીતિ}} (\emph{bhīti}) is a tatsama, Sanskrit \emph{bhīti}, \textquotedblleft{}fear, alarm, dread\textquotedblright{}, as Monier-Williams records it; it shares its root with \textbf{\gu{બીવું}} (\emph{bīvũ}), from the start of the chapter.
\end{cousinweb}

\begin{grammarlens}[title={What you've built}]
One more word for describing what you lived and felt.
\end{grammarlens}

\subsection*{Your turn}
\begin{itemize}
\raggedright
  \item \emph{Say it:} \emph{bhīti}
  \item \emph{Say it:} \emph{bhīti}, once more
  \item \emph{Say it:} \emph{bhūkamp pachī gāmmā̃ ghaṇī bhīti hatī.} — After the earthquake there was great apprehension in the village.
  \item \emph{Recall:} say the Gujarati for to shudder, then the Gujarati for torment, then say \emph{bhīti} again
\end{itemize}

\begin{tcolorbox}[breakable,colback=teal!4,colframe=teal!35!black,title={Before you move on}]
What does \gu{ભીતિ} mean? (\textquotedblleft{}Apprehension; fear of what may come\textquotedblright{}.) Say it once more.
\end{tcolorbox}
